\section{Methodology}

The OCR system is organized as a linear pipeline designed with modularity as a core principle. Each of the six stages has a well-defined interface: it accepts input data, processes it deterministically, and returns output that feeds into other stage. This sequential organization reflects the logical flow of the OCR problem: we must first separate text from background (binarization), then locate where text regions are (row segmentation), then locate individual characters within those regions (character segmentation), then normalize those characters to a standard form (preprocessing), then establish what each character looks like (alphabet creation), and finally match observed characters against known templates (recognition).\\
The six stages are implemented as separate MATLAB functions, each with a clear purpose.\\
\begin{itemize}
    \item Task 1 (\texttt{task1BinarizationAdvanced.m}) converts grayscale or RGB images into binary matrices.
    \item Task 2 (\texttt{task2ExtractRows.m}, \texttt{task2RowProjection.m}, \texttt{task2SegmentRows.m}) locates row boundaries.
    \item Task 3 (\texttt{task3SegmentCharacters.m}) locates character boundaries within rows.
    \item Task 4 (\texttt{task4PreprocessingCharacters.m}) resizes character images to 32×32 pixels.
    \item Task 5 (\texttt{task5CreatingAlphabet.m}) extracts and saves reference characters.
    \item Task 6 (\texttt{task6RecognizeCharacters.m}) matches test characters against alphabet templates.
\end{itemize}